\documentclass[a4paper,11pt]{article}

% ── Encoding & Language ───────────────────────────────────────────────────────
\usepackage[utf8]{inputenc}
\usepackage[T1]{fontenc}
\usepackage[french]{babel}

% ── Layout ────────────────────────────────────────────────────────────────────
\usepackage[top=2.5cm, bottom=2.5cm, left=2.5cm, right=2.5cm]{geometry}
\usepackage{parskip}
\usepackage{setspace}

% ── Typography ────────────────────────────────────────────────────────────────
\usepackage{lmodern}
\usepackage{microtype}

% ── Colours ───────────────────────────────────────────────────────────────────
\usepackage{xcolor}
\definecolor{mainblue}{HTML}{1a2c42}
\definecolor{accentblue}{HTML}{2980b9}
\definecolor{lightblue}{HTML}{d6eaf8}
\definecolor{done}{HTML}{1e8449}
\definecolor{inprog}{HTML}{d35400}
\definecolor{nottodo}{HTML}{922b21}
\definecolor{rowdone}{HTML}{d5f5e3}
\definecolor{rowinprog}{HTML}{fdebd0}
\definecolor{rownot}{HTML}{fadbd8}
\definecolor{rowgray}{HTML}{f0f3f4}
\definecolor{friendly}{HTML}{1e8449}
\definecolor{hostile}{HTML}{922b21}
\definecolor{neutral}{HTML}{626567}

% ── Tables ────────────────────────────────────────────────────────────────────
\usepackage{booktabs}
\usepackage{colortbl}
\usepackage{tabularx}
\usepackage{array}
\usepackage{longtable}
\usepackage{multirow}

% ── Graphics & Diagrams ───────────────────────────────────────────────────────
\usepackage{tikz}
\usetikzlibrary{shapes.geometric, arrows.meta, positioning, fit, backgrounds}
\usepackage{float}

% ── Header / Footer ───────────────────────────────────────────────────────────
\usepackage{fancyhdr}
\pagestyle{fancy}
\fancyhf{}
\fancyhead[L]{\small\textcolor{mainblue}{\textit{Rapport d'avancement S2 — Réseaux de personnages}}}
\fancyhead[R]{\small\textcolor{mainblue}{Mars 2026}}
\fancyfoot[C]{\small\thepage}
\renewcommand{\headrulewidth}{0.4pt}

% ── Sections styling ─────────────────────────────────────────────────────────
\usepackage{titlesec}
\titleformat{\section}{\large\bfseries\color{mainblue}}{}{0em}{\thesection\quad}[\titlerule]
\titleformat{\subsection}{\normalsize\bfseries\color{accentblue}}{}{0em}{\thesubsection\quad}
\titleformat{\subsubsection}{\normalsize\bfseries}{}{0em}{\thesubsubsection\quad}

% ── Hyperlinks ────────────────────────────────────────────────────────────────
\usepackage[colorlinks=true, linkcolor=mainblue, urlcolor=accentblue,
            citecolor=mainblue, pdfauthor={Lotfi Abdallah, Wael Mansoura},
            pdftitle={Rapport d'avancement S2}]{hyperref}

% ── Misc ──────────────────────────────────────────────────────────────────────
\usepackage{enumitem}
\usepackage{mdframed}

% ── Helpers ───────────────────────────────────────────────────────────────────
\newcommand{\done}{\textcolor{done}{\textbf{Terminé}}}
\newcommand{\inprog}{\textcolor{inprog}{\textbf{En cours}}}
\newcommand{\notdone}{\textcolor{nottodo}{\textbf{Non commencé}}}
\newcommand{\high}{\textcolor{red!70!black}{\textbf{Haute}}}
\newcommand{\med}{\textcolor{orange!80!black}{\textbf{Moyenne}}}
\newcommand{\low}{\textcolor{gray}{\textbf{Basse}}}

% ─────────────────────────────────────────────────────────────────────────────
\begin{document}

% ── Page de titre ─────────────────────────────────────────────────────────────
\begin{titlepage}
  \centering
  \vspace*{1.5cm}

  {\Large\bfseries\color{mainblue} Université d'Avignon}\\[0.3cm]
  {\normalsize Master 1 Informatique --- Parcours ILSEN}\\[0.3cm]
  {\normalsize Année universitaire 2025--2026}

  \vspace{1.5cm}
  \noindent\rule{\textwidth}{1.5pt}\\[0.5cm]
  {\LARGE\bfseries\color{mainblue}
    Extraction de réseaux de personnages\\[8pt]
    à partir de textes narratifs}\\[0.4cm]
  {\Large Rapport d'avancement --- Semestre 2}
  \\[0.3cm]
  \noindent\rule{\textwidth}{1.5pt}

  \vspace{1.5cm}

  \begin{minipage}{0.48\textwidth}
    \centering
    \textbf{Étudiants}\\[4pt]
    Lotfi \textsc{Abdallah}\\
    Wael \textsc{Mansoura}
  \end{minipage}
  \hfill
  \begin{minipage}{0.48\textwidth}
    \centering
    \textbf{Date}\\[4pt]
    Mars 2026
  \end{minipage}

  \vspace{1.5cm}

  \begin{mdframed}[backgroundcolor=lightblue, linecolor=accentblue, linewidth=1pt,
                   innertopmargin=10pt, innerbottommargin=10pt,
                   innerleftmargin=14pt, innerrightmargin=14pt]
    \small
    \textbf{Résumé.} Ce rapport présente l'état d'avancement du projet au terme du semestre 2.
    Le prototype fonctionnel livre le pipeline complet du semestre 1 (extraction NER, regroupement d'alias,
    co-occurrences, graphes NetworkX, export GraphML/Kaggle) enrichi du module principal du semestre~2 :
    l'étiquetage automatique des relations entre personnages (\textit{friendly} / \textit{hostile} / \textit{neutral})
    par classification zero-shot via Natural Language Inference (NLI, modèle mDeBERTa multilingue).
    Le rapport détaille également la campagne d'expérimentation prévue en avril et l'anticipation du rapport
    final (article en anglais, rendu \textbf{>~12 mai 2026}).
  \end{mdframed}

  \vfill
  {\small Dépôt Moodle --- AMS Projet 2 --- Séance d'évaluation 1 (10 mars 2026)}
\end{titlepage}

\tableofcontents
\newpage

% =============================================================================
\section{Rappel du contexte}
% =============================================================================

\subsection{Projet}

Le projet consiste à extraire automatiquement des \textbf{réseaux de personnages} à partir de deux
romans d'Isaac Asimov traduits en français :

\begin{itemize}[noitemsep]
  \item \textit{Prélude à Fondation} (19 chapitres : paf0--paf18)
  \item \textit{Les Cavernes d'Acier} (18 chapitres : lca0--lca17)
\end{itemize}

\textbf{Total : 37 chapitres traités.}

Le pipeline extrait les entités nommées (personnages), regroupe leurs alias, détecte les co-occurrences
par fenêtre glissante, et construit un graphe de personnages au format GraphML pour chaque chapitre.
Les résultats sont soumis dans le cadre d'une compétition Kaggle.

\subsection{Bilan du semestre 1 (terminé)}

Le S1 a livré un pipeline fonctionnel de bout en bout :

\begin{itemize}[noitemsep]
  \item Extraction NER multi-modèles (spaCy + Stanza, vote majoritaire)
  \item Filtrage des personnages (LP) et des lieux (LL) avec anti-dictionnaire (\textasciitilde 868 entrées)
  \item Regroupement d'alias (Union-Find + similarité floue + overrides manuels)
  \item Détection des co-occurrences (fenêtre glissante configurable, \texttt{distance\_max}=150)
  \item Génération de graphes NetworkX + export GraphML/CSV
  \item Soumission Kaggle fonctionnelle (37 chapitres)
\end{itemize}

\subsection{Objectif du semestre 2}

\textbf{Étiquetage des relations entre personnages} : enrichir chaque arête du graphe avec un type
de relation (\textit{friendly}, \textit{hostile}, \textit{neutral}) afin de qualifier la nature des
interactions au-delà de la simple co-occurrence.

% =============================================================================
\section{État d'avancement --- Tableau de suivi des exigences}
% =============================================================================

\subsection{Exigences du cahier des charges (S1)}

\setlength{\tabcolsep}{5pt}
\renewcommand{\arraystretch}{1.25}

\begin{table}[H]
  \centering
  \footnotesize
  \begin{tabularx}{\textwidth}{>{\centering}p{0.4cm}
                               X
                               >{\centering}p{1.5cm}
                               >{\centering}p{2.0cm}
                               >{\centering\arraybackslash}p{1.8cm}}
    \toprule
    \textbf{\#} & \textbf{Exigence} & \textbf{Priorité} & \textbf{Responsable} & \textbf{Statut} \\
    \midrule
    \rowcolor{rowdone}1 & Extraction NER (multi-modèles spaCy + Stanza) & \high & Lotfi & \done \\
    \rowcolor{rowdone}2 & Filtrage des personnes (LP + anti-dict.\ \textasciitilde{}868 entrées) & \high & Lotfi & \done \\
    \rowcolor{rowdone}3 & Filtrage des lieux (LL + gazetteer Asimov) & \med & Lotfi & \done \\
    \rowcolor{rowdone}4 & Regroupement des alias (Union-Find + fuzzy + overrides) & \high & Lotfi & \done \\
    \rowcolor{rowdone}5 & Détection des co-occurrences (fenêtre glissante) & \high & Wael & \done \\
    \rowcolor{rowdone}6 & Génération du graphe NetworkX (GraphML) & \high & Wael & \done \\
    \rowcolor{rowdone}7 & Visualisation interactive HTML (PyVis/vis.js) & \med & Wael & \done \\
    \rowcolor{rowdone}8 & Export CSV Kaggle (37 chapitres, GraphML complet) & \high & Lotfi + Wael & \done \\
    \bottomrule
  \end{tabularx}
  \caption{Suivi des exigences S1 --- \textcolor{done}{\textbf{8/8 terminées}}}
\end{table}

\subsection{Exigences supplémentaires du semestre 2}

\begin{table}[H]
  \centering
  \footnotesize
  \begin{tabularx}{\textwidth}{>{\centering}p{0.4cm}
                               X
                               >{\centering}p{1.5cm}
                               >{\centering}p{2.0cm}
                               >{\centering\arraybackslash}p{2.2cm}}
    \toprule
    \textbf{\#} & \textbf{Exigence} & \textbf{Priorité} & \textbf{Responsable} & \textbf{Statut} \\
    \midrule
    \rowcolor{rowdone}9  & Étiquetage des relations (NLI zero-shot) & \high & Lotfi + Wael & \done \\
    \rowcolor{rowdone}10 & Cache des labels de relations (pickle) & \high & Lotfi + Wael & \done \\
    \rowcolor{rowdone}11 & Visualisation colorée des arêtes + HTML combiné & \med & Wael & \done \\
    \rowcolor{rowdone}12 & Rapport de validation (\texttt{print\_validation\_report}) & \med & Lotfi + Wael & \done \\
    \rowcolor{rowdone}13 & NER rule-based sans modèles IA & \low & Lotfi & \done \\
    \rowcolor{rowinprog}14 & Campagne d'expérimentation (avril) & \high & Lotfi + Wael & \inprog \\
    \rowcolor{rownot}15  & Rapport final (article anglais, \LaTeX, > 12 mai) & \high & Lotfi + Wael & \notdone \\
    \bottomrule
  \end{tabularx}
  \caption{Suivi des exigences S2 ---
    \textcolor{done}{\textbf{5 terminées}} \textbullet{}
    \textcolor{inprog}{\textbf{1 en cours}} \textbullet{}
    \textcolor{nottodo}{\textbf{1 non commencée (planifiée)}}}
\end{table}

% =============================================================================
\section{Travaux réalisés au semestre 2}
% =============================================================================

\subsection{Pipeline complet S1 + S2}

La figure~\ref{fig:pipeline} présente le pipeline de bout en bout intégrant les contributions du S1 et
la nouvelle étape d'étiquetage des relations (S2).

\begin{figure}[H]
  \centering
  \begin{tikzpicture}[
    node distance = 0.55cm and 0.2cm,
    box/.style = {rectangle, draw=mainblue, thick, fill=lightblue,
                  rounded corners=4pt, text width=4.2cm,
                  align=center, font=\small, minimum height=0.65cm},
    newbox/.style = {rectangle, draw=done, very thick, fill=rowdone,
                     rounded corners=4pt, text width=4.8cm,
                     align=center, font=\small, minimum height=0.65cm},
    arr/.style = {-Stealth, thick, mainblue},
    lbl/.style = {font=\scriptsize\itshape, text=gray, anchor=west}
  ]
    \node[box]    (ner)   {\textbf{NER}\\rule-based / multi-modèles};
    \node[box, below=of ner]    (filt)  {\textbf{Filtrage LP/LL}\\anti-dict + gazetteer};
    \node[box, below=of filt]   (alias) {\textbf{Regroupement alias}\\Union-Find + fuzzy};
    \node[box, below=of alias]  (cooc)  {\textbf{Co-occurrences}\\fenêtre glissante};
    \node[newbox, below=of cooc](rel)   {\textbf{[S2] Étiquetage relations}\\NLI zero-shot (mDeBERTa)};
    \node[box, below=of rel]    (graph) {\textbf{Graphe NetworkX}\\attribut \texttt{edge\_type}};
    \node[box, below=of graph]  (out)   {\textbf{Sortie}\\GraphML / CSV Kaggle / HTML};

    \draw[arr] (ner)   -- (filt);
    \draw[arr] (filt)  -- (alias);
    \draw[arr] (alias) -- (cooc);
    \draw[arr] (cooc)  -- (rel);
    \draw[arr] (rel)   -- (graph);
    \draw[arr] (graph) -- (out);

    \node[lbl] at ([xshift=0.3cm]ner.east)   {\texttt{nlp\_manual\_ner.py}};
    \node[lbl] at ([xshift=0.3cm]filt.east)  {\texttt{nlp\_extract\_characters.py}};
    \node[lbl] at ([xshift=0.3cm]alias.east) {\texttt{nlp\_aliases.py}};
    \node[lbl] at ([xshift=0.3cm]cooc.east)  {\texttt{nlp\_cooccurrence.py}};
    \node[lbl] at ([xshift=0.3cm]rel.east)   {\texttt{nlp\_relation.py} \textbf{--- NOUVEAU S2}};
    \node[lbl] at ([xshift=0.3cm]graph.east) {\texttt{nlp\_graph.py} \textbf{--- modifié}};
    \node[lbl] at ([xshift=0.3cm]out.east)   {\texttt{nlp\_visualize\_web.py} \textbf{--- modifié}};

    \begin{scope}[on background layer]
      \node[draw=done, fill=rowdone, opacity=0.20, rounded corners=6pt,
            fit=(rel)(graph)(out), inner sep=6pt] {};
    \end{scope}
  \end{tikzpicture}
  \caption{Pipeline complet --- les blocs en vert représentent les contributions du S2}
  \label{fig:pipeline}
\end{figure}

\subsection{Étiquetage des relations — \texttt{nlp\_relation.py} (nouveau)}

\textbf{Objectif :} pour chaque paire de personnages co-occurrents dans un chapitre, attribuer un
label de relation parmi \textcolor{friendly}{\textbf{friendly}}, \textcolor{hostile}{\textbf{hostile}}
ou \textcolor{neutral}{\textbf{neutral}}.

\textbf{Approche retenue : classification zero-shot par NLI.}

\begin{enumerate}[noitemsep]
  \item \textbf{Extraction de contextes} : jusqu'à 5 extraits textuels de 400 caractères entourant chaque
        co-occurrence de la paire (charA, charB).
  \item \textbf{Inférence NLI} : pour chaque extrait, le modèle score trois hypothèses en français :
        « ces deux personnages ont une relation amicale / hostile / neutre ».
  \item \textbf{Vote pondéré} : score moyen par classe sur les 5 extraits.
  \item \textbf{Seuil de confiance} : si le score maximal est inférieur à 0{,}55, le label est
        \textit{neutral} par défaut.
\end{enumerate}

\subsubsection*{Modèle NLI retenu}

\begin{table}[H]
  \centering
  \footnotesize
  \begin{tabular}{lp{9cm}}
    \toprule
    \textbf{Critère} & \textbf{Justification} \\
    \midrule
    Pas de données annotées & Fonctionne en zero-shot, aucun entraînement requis \\
    Support du français     & Entraîné sur XNLI multilingue ; performant en français \\
    Compatible CPU          & \textasciitilde{}1--3\,s/extrait sur Colab gratuit \\
    Déterministe            & Inférence greedy + cache pickle = résultats reproductibles \\
    Simple d'utilisation    & 5 lignes via HuggingFace \texttt{pipeline()} \\
    \bottomrule
  \end{tabular}
  \caption{Justification du choix de \texttt{mDeBERTa-v3-base-mnli-xnli} (\textasciitilde{}350\,Mo)}
\end{table}

\subsubsection*{Approches rejetées}

\begin{table}[H]
  \centering
  \footnotesize
  \begin{tabular}{lp{9.5cm}}
    \toprule
    \textbf{Approche} & \textbf{Raison du rejet} \\
    \midrule
    Analyse de sentiment   & Sentiment $\neq$ relation ; prose expositoire d'Asimov $\Rightarrow$ faux neutres \\
    Règles syntaxiques     & Précision correcte mais rappel quasi nul sur texte littéraire \\
    Fine-tuning Transformer & Nécessite des données annotées que nous n'avons pas \\
    LLM local (GGUF)       & Trop lent sur CPU Colab ; réservé comme voie d'escalade \\
    \bottomrule
  \end{tabular}
  \caption{Approches envisagées et raisons de rejet}
\end{table}

\subsubsection*{Stratégie de cache}

\begin{itemize}[noitemsep]
  \item Clé : \texttt{(chapter\_id, canonA, canonB)} avec canonA $<$ canonB (tri lexicographique)
  \item Fichier : \texttt{relation\_cache.pkl} (pickle Python)
  \item Invalidation : uniquement si changement du modèle NLI ou du seuil de confiance
\end{itemize}

\subsection{NER rule-based — \texttt{nlp\_manual\_ner.py} (nouveau)}

Alternative au NER multi-modèles du S1, fonctionnant \textbf{sans aucun modèle IA} :

\begin{enumerate}[noitemsep]
  \item \textbf{Gazetteer} : dictionnaires des personnages et lieux des deux romans (correspondance la
        plus longue d'abord, non chevauchante).
  \item \textbf{Heuristiques de capitalisation} : mots capitalisés en milieu de phrase, absents d'une
        liste noire de \textasciitilde{}200 mots courants français.
  \item \textbf{Patterns regex} : \textit{Docteur X}, \textit{R. Daneel} (robots asimoviens), préfixes
        de titres (\textit{Professeur}, \textit{Inspecteur}\ldots).
\end{enumerate}

\textbf{Avantage :} exécution quasi-instantanée (quelques secondes pour les 37 chapitres), aucun
téléchargement de modèle, idéal pour le développement itératif rapide.

\subsection{Modifications de \texttt{nlp\_graph.py}}

\begin{itemize}[noitemsep]
  \item Ajout du paramètre \texttt{edge\_labels} à \texttt{generate\_graph()}
  \item Écriture de l'attribut \texttt{edge\_type} sur chaque arête du graphe NetworkX
  \item Le champ \texttt{edge\_type} apparaît automatiquement dans la sortie GraphML
  \item Ajout du reverse alias map pour l'attribut \texttt{names} des nœuds
\end{itemize}

\subsection{Visualisation enrichie — \texttt{nlp\_visualize\_web.py}}

Arêtes colorées par type de relation dans la visualisation HTML :

\begin{itemize}[noitemsep]
  \item \textcolor{friendly}{\rule{9pt}{9pt}}\ \textcolor{friendly}{\textbf{Vert}} $\to$ friendly
  \item \textcolor{hostile}{\rule{9pt}{9pt}}\ \textcolor{hostile}{\textbf{Rouge}} $\to$ hostile
  \item \textcolor{neutral}{\rule{9pt}{9pt}}\ \textcolor{neutral}{\textbf{Gris}} $\to$ neutral
\end{itemize}

Nouvelle fonction \texttt{create\_combined\_html()} : génère un \textbf{seul fichier
\texttt{all\_networks.html}} contenant les 37 chapitres avec :
barre de navigation par livre/chapitre, statistiques (nœuds, arêtes, labels), infobulles au survol
(poids + type de relation), légende visuelle.

\subsection{Notebook intégrateur — \texttt{nomodels.ipynb}}

Notebook complet orchestrant le pipeline de bout en bout avec la NER rule-based :

\begin{table}[H]
  \centering
  \footnotesize
  \begin{tabular}{llr}
    \toprule
    \textbf{Section} & \textbf{Contenu} & \textbf{Durée} \\
    \midrule
    1. Setup              & Imports, rechargement des modules                       & instantané \\
    2. Configuration      & \texttt{books\_config}, anti-dictionnaire               & instantané \\
    3. Extraction NER     & NER rule-based sur 37 chapitres                         & secondes \\
    3.1 Cache NER         & Sauvegarde/chargement pickle (Google Drive)             & instantané \\
    4. Relation cache     & Chargement \texttt{relation\_cache.pkl}                 & instantané \\
    4.1 Boucle principale & Co-occurrences + \texttt{label\_relationships()} + graphe + CSV & heures \\
    4.2 Cache relations   & Sauvegarde du cache mis à jour                          & instantané \\
    5. Vérification       & Aperçu soumission + \texttt{print\_validation\_report()} & instantané \\
    6. Test fenêtres      & Comparaison \texttt{distance\_max} = [25, 50, 75, 100, 150] & secondes \\
    7. Visualisation HTML & Fichier HTML combiné 37 chapitres interactif            & secondes \\
    \bottomrule
  \end{tabular}
  \caption{Sections du notebook \texttt{nomodels.ipynb}}
\end{table}

% =============================================================================
\section{Indicateurs de performance (KPIs)}
% =============================================================================

\subsection{Couverture du corpus}

\begin{table}[H]
  \centering
  \small
  \begin{tabular}{lr}
    \toprule
    \textbf{Indicateur} & \textbf{Valeur} \\
    \midrule
    Chapitres traités         & \textbf{37 / 37} (100\,\%) \\
    Livres couverts           & 2 / 2 (\textit{Prélude à Fondation} + \textit{Les Cavernes d'Acier}) \\
    Nœuds typiques/chapitre   & 5--20 \\
    Arêtes typiques/chapitre  & 5--30 \\
    \bottomrule
  \end{tabular}
  \caption{Couverture du corpus}
\end{table}

\subsection{Paramètres du système}

\begin{table}[H]
  \centering
  \small
  \begin{tabular}{lrl}
    \toprule
    \textbf{Paramètre} & \textbf{Valeur} & \textbf{Rôle} \\
    \midrule
    \texttt{distance\_max}          & 150 mots & Taille de la fenêtre de co-occurrence \\
    \texttt{min\_occurrences}       & 2        & Seuil minimal de mentions d'un personnage \\
    \texttt{CONFIDENCE\_THRESHOLD}  & 0{,}55   & Seuil NLI (en dessous $\Rightarrow$ neutral) \\
    \texttt{MAX\_CONTEXTS\_PER\_PAIR} & 5      & Nombre max d'extraits analysés par paire \\
    \texttt{MAX\_SNIPPET\_CHARS}    & 400      & Longueur max des extraits du modèle NLI \\
    \bottomrule
  \end{tabular}
  \caption{Paramètres configurables du système}
\end{table}

\subsection{Modèle NLI et performances}

\begin{table}[H]
  \centering
  \small
  \begin{tabular}{ll}
    \toprule
    \textbf{Indicateur} & \textbf{Valeur} \\
    \midrule
    Modèle NLI              & \texttt{MoritzLaurer/mDeBERTa-v3-base-mnli-xnli} \\
    Taille du modèle        & \textasciitilde{}350\,Mo \\
    Taxonomie               & 3 classes : friendly / hostile / neutral \\
    Exécution               & CPU uniquement (Google Colab gratuit) \\
    Temps par extrait       & 1--3\,s \\
    NER rule-based (37 chap.) & quelques secondes \\
    Boucle NLI complète     & quelques heures (Colab CPU) \\
    Reproductibilité        & inférence greedy + cache pickle \\
    \bottomrule
  \end{tabular}
  \caption{Performance technique}
\end{table}

\textbf{Note sur la distribution des labels.} Les chapitres d'Asimov sont souvent expositoires ou
philosophiques — de longs passages où les personnages débattent sans vocabulaire émotif. Une proportion
élevée de labels \textit{neutral} est attendue ; ce n'est pas un défaut du système, mais le reflet
du contenu textuel de chaque chapitre.

% =============================================================================
\section{Préparation de la campagne d'expérimentation (avril)}
% =============================================================================

La campagne d'expérimentation (avril 2026) vise à évaluer rigoureusement la qualité de l'étiquetage
des relations et à explorer des axes d'amélioration. Cinq expériences sont planifiées.

\subsection*{Exp.\ 1 — Annotation manuelle (gold standard)}

\textbf{Objectif :} construire un jeu de référence pour l'évaluation quantitative.

\textbf{Protocole :}
\begin{itemize}[noitemsep]
  \item Sélectionner \textasciitilde{}50 paires de personnages (25 par livre),
        représentatives des trois classes.
  \item Chaque membre du binôme annote indépendamment le type de relation en lisant le texte source.
  \item Calcul du taux d'accord inter-annotateurs (Cohen's $\kappa$).
\end{itemize}

\textbf{Livrable :} \texttt{gold\_annotations.csv} (colonnes : \texttt{chapter\_id, charA, charB, gold\_label}).

\subsection*{Exp.\ 2 — Comparaison de modèles NLI}

\textbf{Objectif :} determiner si un modèle alternatif produit des labels de meilleure qualité.

\begin{table}[H]
  \centering
  \small
  \begin{tabular}{lp{8cm}}
    \toprule
    \textbf{Modèle} & \textbf{Caractéristique} \\
    \midrule
    \texttt{mDeBERTa-v3-base-mnli-xnli}     & Multilingue, baseline actuelle \\
    \texttt{cmarkea/distilcamembert-base-nli} & Spécifique au français, plus léger \\
    Mistral API (\texttt{mistral-medium}, T=0) & LLM en batch, qualité quasi-humaine \\
    \bottomrule
  \end{tabular}
  \caption{Modèles comparés dans l'Exp.\ 2}
\end{table}

\textbf{Protocole :} évaluer chaque modèle sur le gold standard (précision, rappel, F1 macro).

\subsection*{Exp.\ 3 — Impact de \texttt{distance\_max} sur la qualité}

\textbf{Objectif :} évaluer si la taille de la fenêtre affecte la qualité des labels (et pas seulement
le nombre de co-occurrences).

\textbf{Protocole :} tester \texttt{distance\_max} $\in$ [25, 50, 75, 100, 150, 200] ;
comparer la distribution des labels et la qualité sur le gold standard.

\subsection*{Exp.\ 4 — Faisabilité d'une taxonomie 5 classes}

\textbf{Objectif :} explorer l'extension à 5 classes.

\begin{table}[H]
  \centering
  \small
  \begin{tabular}{ll}
    \toprule
    \textbf{Label} & \textbf{Signification} \\
    \midrule
    \texttt{ally}      & Coopération avec un objectif commun \\
    \texttt{mentor}    & Relation de guidance/enseignement \\
    \texttt{hostile}   & Conflit, opposition \\
    \texttt{romantic}  & Relation intime \\
    \texttt{neutral}   & Pas de signal clair \\
    \bottomrule
  \end{tabular}
  \caption{Taxonomie étendue à 5 classes}
\end{table}

\subsection*{Exp.\ 5 — LLM comme baseline de qualité supérieure}

\textbf{Objectif :} évaluer si un appel batch à un LLM peut servir de référence ou de remplacement
pour le modèle NLI.

\textbf{Protocole :} appel Mistral API avec \texttt{temperature=0} :
\begin{quote}
  \textit{« Dans ce passage, la relation entre \{charA\} et \{charB\} est-elle amicale, hostile
  ou neutre ? Réponds par un seul mot. »}
\end{quote}
Comparaison des résultats avec les labels NLI sur le gold standard.

\subsection{Calendrier de la campagne d'expérimentation}

\begin{table}[H]
  \centering
  \small
  \begin{tabular}{ll}
    \toprule
    \textbf{Semaine} & \textbf{Activité} \\
    \midrule
    Semaine 1 (avril) & Exp.\ 1 : annotation manuelle, gold standard (50 paires) \\
    Semaine 2 (avril) & Exp.\ 2 : comparaison modèles NLI \\
    Semaine 3 (avril) & Exp.\ 3 : impact de \texttt{distance\_max} \\
    Semaine 4 (avril) & Exp.\ 4 \& 5 : taxonomie 5 classes + LLM baseline \\
    \bottomrule
  \end{tabular}
  \caption{Calendrier de la campagne d'expérimentation}
\end{table}

% =============================================================================
\section{Anticipation du rapport final (> 12 mai 2026)}
% =============================================================================

\subsection{Format imposé}

Le rapport final est un \textbf{article scientifique en anglais}, rédigé en \LaTeX{}, avec la
structure suivante :

\begin{enumerate}[noitemsep]
  \item \textbf{Introduction} --- contexte, problématique, revue de littérature
  \item \textbf{Methodology} --- pipeline complet, choix techniques justifiés
  \item \textbf{Experimental Results} --- résultats de la campagne d'expérimentation (avril)
  \item \textbf{Conclusion} --- bilan, limites, perspectives
  \item \textbf{References} --- bibliographie complète
\end{enumerate}

\subsection{Bibliographie déjà identifiée}

\begin{table}[H]
  \centering
  \footnotesize
  \begin{tabular}{lp{10cm}}
    \toprule
    \textbf{Thème} & \textbf{Référence} \\
    \midrule
    Zero-shot NLI         & Yin et al.\ (2019) --- \textit{Benchmarking Zero-shot Text Classification} \\
    mDeBERTa / DeBERTa    & He et al.\ (2021) --- \textit{DeBERTa: Decoding-enhanced BERT with Disentangled Attention} \\
    XNLI (multilingual)   & Conneau et al.\ (2018) --- \textit{XNLI: Evaluating Cross-lingual Sentence Representations} \\
    Réseaux de personnages & Labatut \& Bost (2019) --- \textit{Extraction and Analysis of Fictional Character Networks} \\
    Réseaux de personnages & Elson \& McKeown (2010) --- \textit{Automatic Attribution of Quoted Speech in Literary Narrative} \\
    Résolution d'alias    & Vala et al.\ (2015) --- \textit{Mr.\ Bennet, his coachman, and the Archbishop walk into a bar\ldots} \\
    NER en français       & spaCy / Stanza documentation et articles de référence \\
    \bottomrule
  \end{tabular}
  \caption{Bibliographie cible pour le rapport final}
\end{table}

\subsection{Planning de rédaction}

\begin{table}[H]
  \centering
  \small
  \begin{tabular}{ll}
    \toprule
    \textbf{Période} & \textbf{Activité} \\
    \midrule
    Mi-avril          & Collecte des résultats expérimentaux + bibliographie complète \\
    Fin avril         & Rédaction Methodology + Experimental Results \\
    Semaine 1 de mai  & Rédaction Introduction + revue de littérature \\
    Semaine 2 de mai  & Conclusion + relecture croisée \\
    \textbf{12 mai}   & \textbf{Soumission du rapport final} \\
    \bottomrule
  \end{tabular}
  \caption{Planning de rédaction du rapport final}
\end{table}

% =============================================================================
\section{Répartition des tâches}
% =============================================================================

\subsection{Travaux réalisés (S2)}

\begin{table}[H]
  \centering
  \small
  \begin{tabular}{lp{2.5cm}}
    \toprule
    \textbf{Tâche} & \textbf{Responsable} \\
    \midrule
    Conception du plan d'implémentation S2                         & Lotfi + Wael \\
    \texttt{nlp\_relation.py} (NLI zero-shot, cache, validation)   & Lotfi + Wael \\
    \texttt{nlp\_manual\_ner.py} (NER rule-based)                  & Lotfi \\
    \texttt{nlp\_graph.py} (attribut \texttt{edge\_type})          & Wael \\
    \texttt{nlp\_visualize\_web.py} (arêtes colorées + HTML combiné) & Lotfi \\
    Notebook \texttt{nomodels.ipynb} (intégration bout en bout)    & Wael \\
    Mise à jour \texttt{antidict.txt} + \texttt{characters.txt}    & Lotfi \\
    \bottomrule
  \end{tabular}
  \caption{Répartition des travaux réalisés au S2}
\end{table}

\subsection{Travaux à venir}

\begin{table}[H]
  \centering
  \small
  \begin{tabular}{lll}
    \toprule
    \textbf{Tâche} & \textbf{Responsable} & \textbf{Échéance} \\
    \midrule
    Annotation manuelle (gold standard, 50 paires) & Lotfi + Wael & Début avril \\
    Campagne d'expérimentation (Exp.\ 1--5)        & Lotfi + Wael & Avril \\
    Collecte des références bibliographiques       & Lotfi        & Mi-avril \\
    Rédaction rapport final (\LaTeX, anglais)      & Lotfi + Wael & Fin avril -- mai \\
    Relecture croisée et soumission                & Lotfi + Wael & 12 mai \\
    \bottomrule
  \end{tabular}
  \caption{Tâches à venir et échéances}
\end{table}

% =============================================================================
\newpage
\appendix
\section{Annexe — Inventaire des fichiers du projet}
% =============================================================================

\begin{table}[H]
  \centering
  \footnotesize
  \begin{tabularx}{\textwidth}{lp{1.6cm}X}
    \toprule
    \textbf{Fichier} & \textbf{Statut S2} & \textbf{Rôle} \\
    \midrule
    \texttt{nlp\_relation.py}           & \textbf{Nouveau}  & Classification des relations (NLI zero-shot, cache, validation) \\
    \texttt{nlp\_manual\_ner.py}        & \textbf{Nouveau}  & NER rule-based sans modèle IA \\
    \texttt{nomodels.ipynb}             & \textbf{Nouveau}  & Notebook complet (pipeline rule-based + étiquetage) \\
    \texttt{nlp\_graph.py}              & \textbf{Modifié}  & Génération du graphe avec attribut \texttt{edge\_type} \\
    \texttt{nlp\_visualize\_web.py}     & \textbf{Modifié}  & Visualisation interactive avec arêtes colorées \\
    \texttt{antidict.txt}               & \textbf{Modifié}  & Anti-dictionnaire (\textasciitilde{}868 entrées) \\
    \texttt{characters.txt}             & \textbf{Modifié}  & Liste de personnages pour le gazetteer \\
    \texttt{nlp\_cooccurrence.py}       & Inchangé          & Détection des co-occurrences \\
    \texttt{nlp\_extract\_characters.py}& Inchangé          & Comptage et filtrage des entités \\
    \texttt{nlp\_aliases.py}            & Inchangé          & Regroupement d'alias (Union-Find) \\
    \texttt{nlp\_utils.py}              & Inchangé          & Utilitaires (lecture de fichiers, anti-dict) \\
    \texttt{all\_networks.html}         & \textbf{Nouveau}  & Visualisation combinée 37 chapitres (arêtes colorées) \\
    \bottomrule
  \end{tabularx}
  \caption{Inventaire des fichiers du projet}
\end{table}

\end{document}
